\subsection{应用：I. 有理整数}

考虑自然数集的交换幺半群 \(\mathbf{N}\) ，其合成律为加法； \(\mathbf{N}\) 的所有元素在此律下都是可消去的（集合论，III，§5，第2号，推论3）。 \(\mathbf{N}\) 的差群记为 \(\mathbf{Z}\) ，其元素称为有理整数；其律称为有理整数加法，也记作 \(+\) 。从 \(\mathbf{N}\) 到 \(\mathbf{Z}\) 的典范同态是单射的，我们将把 \(\mathbf{N}\) 与它在 \(\mathbf{Z}\) 中的像等同。 \(\mathbf{Z}\) 的元素根据定义是由 \(\mathbf{N} \times \mathbf{N}\) 中 \((m_1, n_1)\) 和 \((m_2, n_2)\) 之间写作 \(m_1 + n_2 = m_2 + n_1\) 的关系所确定的等价类；一个元素 \(m\) （ \(\mathbf{N}\) 的）与由元素 \((m + n, n)\) （其中 \(n \in N\) ）组成的类等同；它在 Z 中的负元是元素 \((n, m + n)\) 的类。每个元素 \((p, q)\) （ \(N \times N\) 的）都可以写成 \((m + n, n)\) 的形式（若 \(p \geqslant q\) ），或 \((n, m + n)\) 的形式（若 \(p \leqslant q\) ）；由此可知，Z 是 N 与 N 中元素的负元所成集合的并集。单位元 0 是 N 中唯一一个其负元也属于 N 的元素。

对于每个自然数 m，-m 表示 m 的负有理整数，且 -N 表示所有元素 -m（其中 \(m \in N\) ）所成的集合。则

\[
\mathbf {Z} = \mathbf {N} \cup (- \mathbf {N}) \quad \text { and } \quad \mathbf {N} \cap (- \mathbf {N}) = \{0 \}.
\]

对于 \(m \in \mathbf{N}\) , \(m = -m\) 当且仅当 \(m = 0\) .

设 \(m\) 和 \(n\) 是两个自然数；

\begin{enumerate}
\def\labelenumi{(\alph{enumi})}
\item
  若 \(m \geqslant n\) ，则 \(m + (-n) = p\) ，其中 \(p\) 是 \(\mathbf{N}\) 中使得 \(m = n + p\) 的元素;
\item
  若 \(m \leqslant n\) ，则 \(m + (-n) = -p\) ，其中 \(p\) 是 \(\mathbf{N}\) 中使得 \(m + p = n\) 的元素;
\end{enumerate}

\[
(c) (- m) + (- n) = - (m + n).
\]

性质(b)和(c)由第3号命题4推出；由于 \(\mathbf{Z} = \mathbf{N} \cup (-\mathbf{N})\) ， \(\mathbf{N}\) 中的加法连同性质(a)、(b)和(c)完全描述了 \(\mathbf{Z}\) 中的加法.

更一般地，-x 用来表示任意有理整数 x 的负元；合成 \(x + (-y)\) 简写为 x - y（参见第8号）。

自然数之间的序关系 \(\leqslant\) 由以下性质刻画： \(m \leqslant n\) 当且仅当存在整数 \(p \in N\) 使得 \(m + p = n\) （集合论，III，§3，第6号，命题13及§5，第2号，命题2）。有理整数 x 与 y 之间的关系 \(y - x \in N\) 是 Z 上的一个全序关系，它延拓了 N 上的序关系 \(\leqslant\) 。事实上，对所有 \(x \in Z\) , \(x - x = 0 \in N\) ；若 \(y - x \in N\) 和 \(z - y \in N\) ，则

\[
z - x = (z - y) + (y - x) \in \mathbf {N},
\]

因为 \(\mathbf{N}\) 在加法下是稳定的；若 \(y - x \in \mathbf{N}\) 和 \(x - y \in \mathbf{N}\) ，则 \(y - x = 0\) ，因为 \(0\) 是 \(\mathbf{N}\) 中唯一其负元属于 \(\mathbf{N}\) 的元素；对任意有理整数 \(x\) 和 \(y, y - x \in \mathbf{N}\) 或 \(x - y \in \mathbf{N}\) ，因为 \(\mathbf{Z} = \mathbf{N} \cup (-\mathbf{N})\) ；最后，若 \(x\) 和 \(y\) 是自然数，则 \(y - x \in \mathbf{N}\) 当且仅当存在 \(p \in \mathbf{N}\) 使得 \(x + p = y\) 。此序关系也记作 \(\leqslant\) .

今后，当 Z 被视为有序集时，除非另有说明，它总是采用刚刚定义的序，自然数与整数 \(\geqslant 0\) 视为同一；它们也被称为正整数；整数 \(\leqslant 0\) ，即正整数的负元，被称为负整数；整数 \textgreater{} 0（相应地 \textless{} 0）被称为严格正的（相应地严格负的）；整数 \textgreater{} 0 所成的集合有时记为N*。

设 \(x, y\) 和 \(z\) 是三个有理整数；则 \(x \leqslant y\) 当且仅当 \(x + z \leqslant y + z\) 。因为 \(x - y = (x + z) - (y + z)\) 。这一性质通过说 Z 上的序关系在平移下不变来表达。

